\Needspace{8\baselineskip}\originalchapter{2006 年期中卷}
题源: 31 Oct 2006, 1.5小时、闭卷允许一张8.5\ensuremath{\times}11英寸笔记. 三题权重35\%、35\%、30\%. 当期官方答案未在所得站点文件中发现, 以下独立解答.
\begin{exercise}\label{ex:06-mid-1}
原题1(35\%). 二维偶极场 $\mathbf B=B_0R^2r^{-2}(-\sin\theta\mathbf e_r+\cos\theta\mathbf e_\theta)$, 与z无关, 粒子m/q. 提示场线为r\ensuremath{\propto}cos\ensuremath{\theta}. $\theta=0$处$r=r_0$, W\ensuremath{\perp}0及总W0. (a)最大角位移\ensuremath{\theta}max; (b)漂移仅用\ensuremath{\theta}和固定参数表示; (c)沿场运动 $\ddot\theta=f(\theta)$求f.
\end{exercise}
\begin{solution}
场线r=r0cos\ensuremath{\theta}, 圆弧线元ds=r0d\ensuremath{\theta}, $B=B_c\sec^2\theta$, Bc=B0R\ensuremath{{}^{2}}/r0\ensuremath{{}^{2}}. \ensuremath{\mu}=W\ensuremath{\perp}0/Bc, 转向 $\mu B=W_0$, 故 $\theta_{\max}=\arccos\sqrt{W_{\perp0}/W_0}$. 梯度 $\nabla B=-2B\mathbf e_r/r$, 曲率 $\boldsymbol\kappa=-(2\cos\theta/r)(\cos\theta\mathbf e_r+\sin\theta\mathbf e_\theta)$. 两漂移沿z, 相加
\begin{equation}
 \mathbf v_D=\frac{2\cos\theta}{qBr}(W_\perp+2W_\parallel)\mathbf e_z
 =\frac{2[2W_0\cos^2\theta-W_{\perp0}]}{qB_cr_0}\mathbf e_z.
\end{equation}
沿场力 $mr_0\ddot\theta=-\mu(\dd B/\dd\theta)/r_0$, 因而 $f=-2W_{\perp0}\sin\theta/(mr_0^2\cos^3\theta)$. 尺度小于回旋半径的近原点不适用导引中心, 但有限镜点轨道不需到$r=0$.
\end{solution}
\begin{exercise}\label{ex:06-mid-2}
原题2(35\%). 第1粒子m1,v0沿x, 冲量参数b, 第2初静且先视无限质量; 排斥径向力K/r\ensuremath{{}^{3}}. (a)最近rmin; (b)偏折\ensuremath{\psi}和90°冲量参数b90; (c)改第2质量m2有限, 以\ensuremath{\psi}表示第1的x动量损失; (d)束进入密度n2静靶海, $\nu_p=n_2\sigma_pv_x$, $\sigma_p=\int_0^\infty f(b)\dd b$, 求f. 原卷附可直接推导的基本积分表.
\end{exercise}
\begin{solution}
势能K/(2r\ensuremath{{}^{2}}), 能量和角动量L=\ensuremath{\mu}vb给
$\dot r^2=v^2[1-(b^2+K/(\mu v^2))/r^2]$.
(a)无限质量\ensuremath{\mu}=m1, $r_{\min}=\sqrt{b^2+K/(m_1v_0^2)}$. (b)一半径向角积分 $\int_{rmin}^\infty bv\dd r/(r^2\dot r)=\pi b/(2rmin)$, 所以 $\psi=\pi(1-b/\sqrt{b^2+a_K^2})$, $a_K^2=K/(\mu v_0^2)$, $b_{90}=a_K/\sqrt3$.

(c)有限质量用\ensuremath{\mu}=m1m2/(m1+m2)求\textbf{相对速度}偏折\ensuremath{\psi}. 实验室出vx=$v_0[m_1+m_2\cos\psi]/(m_1+m_2)$, 损失 $\Delta p_{1x}=m_1v_0m_2(1-\cos\psi)/(m_1+m_2)$. 若把\ensuremath{\psi}改指实验室角\ensuremath{\theta}, 需 $\tan\theta=m_2\sin\psi/(m_1+m_2\cos\psi)$变换, 不可沿用重靶关系.

(d)每碰撞的有向动量损失比例乘环截面,
\begin{equation}
 f(b)=\frac{2\pi b m_2}{m_1+m_2}\left[1-\cos\left(\pi\left(1-\frac b{\sqrt{b^2+K/(\mu v_x^2)}}\right)\right)\right].
\end{equation}
远场\ensuremath{\psi}\ensuremath{\propto}b$^{-2}$, f\ensuremath{\propto}b$^{-3}$, 积分收敛, 没有Coulomb 1/r\ensuremath{{}^{2}}力的对数尾. 原件的积分提示可由变量r/rmin得到, 不依赖第三方积分表.
\end{solution}
\begin{exercise}\label{ex:06-mid-3}
原题3(30\%). 无限z pinch,B\ensuremath{\theta}(r),p(r), 半径a,p(a)=0,\ensuremath{\rho}=\ensuremath{\rho}0常数. 平衡不强制p\ensuremath{\rho}$^{-\gamma}$常数, 压力可独立于密度. (a)刚体转动V\ensuremath{\theta}=r\ensuremath{\Omega}求平衡关系; (b)Jz=J0常数求p(r); (c)\ensuremath{\beta}定义为平均p除 $B_\theta(a)^2/(2\mu_0)$, 求\ensuremath{\beta}.
\end{exercise}
\begin{solution}
向心加速度 $-r\Omega^2\mathbf e_r$, 所以 $p'=\rho_0\Omega^2r-JB$, 或 $(p+B^2/(2\mu_0))'+B^2/(\mu_0r)=\rho_0\Omega^2r$. 安培 $B=\mu_0J_0r/2$, 由边界
$p=(\mu_0J_0^2/4-\rho_0\Omega^2/2)(a^2-r^2)$.
平均抛物线p为中心一半, 得 $\beta=1-2\rho_0\Omega^2/(\mu_0J_0^2)$. 非负压力要求\ensuremath{\rho}\ensuremath{\Omega}\ensuremath{{}^{2}}\ensuremath{\le}$\mu_0$J0\ensuremath{{}^{2}}/2. 转动离心支持使所需内压降低; $\Omega=0$给$\beta=1$, 与静态z pinch积分关系一致.
\end{solution}
